% =============================================================================
% CHAPITRE 6 — CONCLUSION GÉNÉRALE ET PERSPECTIVES
% =============================================================================

\chapter{Conclusion générale et perspectives}

% =============================================================================
% 6.1 CONCLUSION GÉNÉRALE
% =============================================================================

\section{Conclusion générale}

Ce projet de stage a porté sur la conception et la réalisation de Droppy,
une solution intelligente dédiée à la supervision et à l'analyse des données
issues d'un réseau d'eau.

L'objectif principal était de proposer une plateforme capable d'exploiter
les données de télémétrie afin d'identifier automatiquement des comportements
anormaux et de fournir des informations utiles à la supervision du réseau.

La réalisation du projet a permis de mettre en place une architecture
distribuée intégrant une interface de supervision, un service central
d'orchestration et un microservice dédié aux traitements d'intelligence
artificielle. Cette organisation a permis de séparer les différentes
responsabilités tout en assurant leur communication.

La conception de la solution a été formalisée à travers plusieurs modèles
UML permettant de représenter les interactions, l'architecture, les
séquences principales, la structure des classes et l'organisation du
déploiement.

La réalisation de l'application a ensuite permis de concrétiser ces choix
à travers un dashboard destiné à la supervision des anomalies, à la
visualisation des données et à la consultation des prévisions.

La composante intelligente constitue également un élément essentiel de
Droppy. L'approche retenue combine plusieurs mécanismes complémentaires,
notamment des règles métier, l'Isolation Forest et la prévision de
consommation. Cette combinaison permet d'obtenir une analyse plus adaptée
aux caractéristiques des données étudiées.

Ainsi, ce projet a permis de mettre en pratique différentes compétences
liées au développement logiciel, à la conception d'architectures
distribuées, au traitement des données et à l'intelligence artificielle,
tout en répondant à une problématique concrète de supervision.


% =============================================================================
% 6.2 APPORTS DU PROJET
% =============================================================================

\section{Apports du projet}

La réalisation de Droppy a permis d'aboutir à une solution intégrant
plusieurs fonctionnalités complémentaires :

\begin{itemize}

    \item supervision des anomalies détectées ;

    \item visualisation des informations à travers un dashboard ;

    \item analyse automatique des comportements inhabituels ;

    \item génération de prévisions de consommation ;

    \item transmission des événements en temps réel ;

    \item centralisation des résultats pour leur exploitation par
    l'utilisateur.

\end{itemize}

Au-delà des fonctionnalités développées, ce projet a également constitué une
expérience permettant de mieux comprendre les contraintes liées à
l'intégration de plusieurs services logiciels au sein d'une même solution.


% =============================================================================
% 6.3 LIMITES DE LA SOLUTION
% =============================================================================

\section{Limites de la solution}

Malgré les fonctionnalités réalisées, certaines améliorations restent
possibles.

Les performances des mécanismes de détection dépendent notamment de la
qualité et de la disponibilité des données historiques. La diversité des
comportements observés sur différents dispositifs peut également nécessiter
une adaptation des modèles et des paramètres utilisés.

Par ailleurs, l'évolution de la solution vers un environnement de production
à grande échelle nécessiterait de renforcer certains aspects liés au
déploiement, à la supervision des services et à la gestion d'un nombre plus
important de dispositifs.


% =============================================================================
% 6.4 PERSPECTIVES
% =============================================================================

\section{Perspectives}

Plusieurs évolutions peuvent être envisagées afin d'améliorer les capacités
de Droppy.

\subsection{Amélioration des modèles d'intelligence artificielle}

Une première perspective consiste à enrichir les modèles d'analyse en
utilisant davantage de données historiques et en expérimentant différentes
approches d'apprentissage automatique afin d'améliorer la détection des
comportements atypiques.

\subsection{Extension de la supervision}

La solution pourrait être étendue afin de prendre en charge un nombre plus
important de dispositifs et de réseaux. Une gestion plus avancée des
dispositifs permettrait également d'adapter les analyses aux caractéristiques
de chaque installation.

\subsection{Amélioration du système d'alertes}

Le système de supervision pourrait être complété par différents canaux
d'alerte afin de permettre une notification plus rapide des événements
importants.

\subsection{Déploiement dans un environnement Cloud}

Une évolution vers une infrastructure Cloud permettrait de faciliter le
déploiement, la disponibilité et la montée en charge de la solution.

\subsection{Développement d'une application mobile}

Une application mobile pourrait également être développée afin de permettre
aux opérateurs de consulter les anomalies et les alertes depuis leurs
appareils mobiles.

\subsection{Évolution vers une maintenance prédictive}

À plus long terme, les données historiques et les résultats des modèles
pourraient être exploités afin d'aller au-delà de la simple détection
d'anomalies et de proposer une approche de maintenance prédictive.


% =============================================================================
% 6.5 CONCLUSION FINALE
% =============================================================================

\section{Conclusion finale}

Le développement de Droppy a permis de mettre en œuvre une solution complète
associant supervision, analyse de données, intelligence artificielle et
visualisation.

Le projet constitue une première base pour une plateforme pouvant évoluer
vers une supervision plus automatisée et plus prédictive des réseaux d'eau.

Les perspectives identifiées offrent ainsi plusieurs possibilités
d'évolution, aussi bien sur le plan fonctionnel que sur le plan technique.
Elles permettraient notamment d'améliorer la précision des analyses,
d'étendre la solution à davantage de dispositifs et de faciliter son
exploitation dans un environnement réel à grande échelle.